% 601.433/633 Introduction to Algorithms — homework template
% Usage: copy to hw/hwN/hwN.tex, fill in \hwnumber, \duedate, \groupmembers, then
%        tools/build.sh hw/hwN/hwN.tex
% Compiles with the TeX Live packages already on this machine. Pseudocode support
% (algorithm2e / algpseudocode) needs `sudo apt install texlive-science`; see bottom.
\documentclass[11pt,letterpaper]{article}

\usepackage[margin=1in]{geometry}
\usepackage{amsmath,amssymb,amsthm,mathtools}
\usepackage{enumitem}
\usepackage{graphicx}
\usepackage{xcolor}
\usepackage{hyperref}
\hypersetup{colorlinks=true,linkcolor=blue!60!black,urlcolor=blue!60!black}
\usepackage{tikz}
\usepackage{fancyhdr}
\usepackage{setspace}
\usepackage{seqsplit}   % line-breaks the long opening/commitment hex strings

% ---- fill these in -------------------------------------------------------
\newcommand{\hwnumber}{N}
\newcommand{\duedate}{Month DD, 2026}
\newcommand{\studentname}{Carlos Ortiz Montes de Oca}
\newcommand{\groupmembers}{none}        % list collaborators, or "none"
% --------------------------------------------------------------------------

\pagestyle{fancy}
\fancyhf{}
\lhead{601.433/633 Intro to Algorithms}
\chead{Homework \hwnumber}
\rhead{\studentname}
\cfoot{\thepage}

% theorem-like environments (amsthm)
\theoremstyle{plain}
\newtheorem{theorem}{Theorem}
\newtheorem{lemma}{Lemma}
\newtheorem{claim}{Claim}
\newtheorem{corollary}{Corollary}
\theoremstyle{definition}
\newtheorem{definition}{Definition}
\theoremstyle{remark}
\newtheorem*{remark}{Remark}

% problem / part structure
\newcounter{problem}
\newcommand{\problem}[1]{\stepcounter{problem}\section*{Problem \theproblem: #1}}
\newcommand{\parta}[1]{\subsection*{(#1)}}

% handy macros
\DeclarePairedDelimiter\ceil{\lceil}{\rceil}
\DeclarePairedDelimiter\floor{\lfloor}{\rfloor}
\DeclarePairedDelimiter\abs{\lvert}{\rvert}
\DeclarePairedDelimiter\card{\lvert}{\rvert}
\newcommand{\R}{\mathbb{R}}
\newcommand{\Z}{\mathbb{Z}}
\newcommand{\N}{\mathbb{N}}
\newcommand{\E}{\mathbb{E}}
\renewcommand{\Pr}[1]{\mathbb{P}\left[#1\right]}
\DeclareMathOperator{\polylog}{polylog}
\newcommand{\todo}[1]{\textcolor{red}{\textbf{[TODO: #1]}}}   % transcription flags; must be empty before submit
% AI-usage transparency line: \aireport{url}{key}{sha256}. Generated by tools/seal.py at submission time.
\newcommand{\aireport}[3]{\par\noindent\footnotesize\textbf{AI-usage report:} \url{#1} \quad
  \textbf{opening:} \texttt{\seqsplit{#2}} \quad \textbf{Pedersen commitment:} \texttt{\seqsplit{#3}}\normalsize\par}

\setstretch{1.15}   % room for grader comments

\begin{document}

\noindent{\large\textbf{601.433/633 Introduction to Algorithms} \hfill \textbf{Fall 2026}}\\
\textbf{Homework \#\hwnumber} \hfill \textbf{Due:} \duedate\\
\textbf{Name:} \studentname \hfill \textbf{Collaborators:} \groupmembers\\
\rule{\textwidth}{0.4pt}
% \aireport{https://jhu.nymo.finance/hw/N}{OPENING}{COMMITMENT}   <- paste output of tools/seal.py N here

% =========================================================================
\problem{Title (points)}

\parta{a}
\begin{claim}
State exactly what is being proved, as a single sentence.
\end{claim}
\begin{proof}
Proof goes here. Say what you are assuming, what you are proving, and which
technique (direct, contradiction, induction on $n$, potential function, \dots).
\end{proof}

\parta{b}
\todo{transcribe}

% =========================================================================
\problem{Title (points)}

\todo{transcribe}

\end{document}

% --- after `sudo apt install texlive-science`, pseudocode is available: -----
% \usepackage{algorithm}
% \usepackage{algpseudocode}
% \begin{algorithm}[H]
% \caption{Find($x$)}
% \begin{algorithmic}[1]
% \While{$x.\mathit{parent} \neq x$}
%   \State $x \gets x.\mathit{parent}$
% \EndWhile
% \State \Return $x$
% \end{algorithmic}
% \end{algorithm}
